\par\noindent\textbf{(1)} 动能 \(T=\tfrac I2(\dot\theta^2+\dot\phi^2\sin^2\theta)+\tfrac{I_3}{2}(\dot\psi+\dot\phi\cos\theta)^2\)，重力势 \(U=mgl\cos\theta\)。故 \(L=T-U\)，循环动量 \(p_\psi=I_3(\dot\psi+\dot\phi\cos\theta)=s\)，\(p_\phi=I\dot\phi\sin^2\theta+s\cos\theta=j\) 均守恒。
\par\noindent\textbf{(2)} 消去循环速度，约化能量为 \(E=\tfrac I2\dot\theta^2+V_{\mathrm{eff}}(\theta)+s^2/(2I_3)\)，其中 \(V_{\mathrm{eff}}(\theta)=\frac{(j-s\cos\theta)^2}{2I\sin^2\theta}+mgl\cos\theta\)。章动转折角是 \(E-s^2/(2I_3)=V_{\mathrm{eff}}\) 的根，存在性取决于能量和两个守恒动量。
\par\noindent\textbf{(3)} 定常进动取 \(\theta=\theta_0\)、\(\dot\phi=\Omega\) 常数。由 \(\theta\) 的Euler方程并以 \(s=I_3\omega_3\) 表示自旋动量，得 \(I\Omega^2\cos\theta_0-s\Omega+mgl=0\)。故 \(\Omega_\pm=[s\pm\sqrt{s^2-4Imgl\cos\theta_0}]/(2I\cos\theta_0)\)，当 \(\cos\theta_0=0\) 时应直接用原方程而非除零公式。
\par\noindent\textbf{(4)} 若 \(s^2\gg 4Imgl|\cos\theta_0|\)，慢进动 \(\Omega_-\approx mgl/s\)，快进动 \(\Omega_+\approx s/(I\cos\theta_0)\)，两者可分辨。定常倾角还必须满足 \(V_{\mathrm{eff}}'(\theta_0)=0\)，小章动稳定需 \(V_{\mathrm{eff}}''(\theta_0)>0\)，频率平方为 \(V_{\mathrm{eff}}''/I\)，仅判别式为正不够。
